% Part: first-order-logic
% Chapter: syntax-and-semantics
% Section: Structures

\documentclass[../../../include/open-logic-section]{subfiles}

\begin{document}

\olfileid{fol}{syn}{str}

\olsection{\printtoken{P}{structure} kanggo Basa Logika Tataran Kapisan}

\begin{explain}
Basa logika tataran kapisan, yen mung basa iku dhewe, \emph{durung diwenehi
interpretasi:} !!{constant}s, !!{function}s, lan !!{predicate}s durung
nduweni teges tartamtu. Teges diwenehake kanthi nemtokake
\article{structure} \emph{!!{structure}}. Struktur iki nemtokake
\emph{domain}, yaiku obyek-obyek kang dituduhake dening !!{constant}s,
dadi papan tumindake !!{function}s, lan dadi wilayah kang dilewati
dening kuantor. Kajaba iku, struktur nemtokake !!{constant}s endi
kang nuduhake obyek endi, carane !!a{function} masangake obyek
marang obyek, lan obyek endi kang dicakup dening !!{predicate}s.
!!^{structure}s dadi dhasar kanggo gagasan \emph{semantis} ing logika,
umpamane gagasan akibat, validitas, lan bisa dicukupi. Ing pustaka,
struktur uga diarani ``struktur,'' ``interpretasi,'' utawa ``model.''
\end{explain}

\begin{defn}[!!^{structure}s]
\Article{structure} \emph{!!{structure}}~$\Struct M$ kanggo basa
$\Lang{L}$ saka logika tataran kapisan dumadi saka unsur-unsur iki:
\begin{enumerate}
\item \emph{Domain:} himpunan kang ora kosong, $\Domain M$
\item \emph{Interpretasi !!{constant}s:} kanggo saben !!{constant}~$c$
  ing $\Lang{L}$, !!a{element} $\Assign{c}{M} \in \Domain M$
\item \emph{Interpretasi !!{predicate}s:} kanggo saben $n$-panggonan
  !!{predicate}~$R$ ing $\Lang{L}$ (saliyane $\eq$), relasi
  $n$-panggonan $\Assign{R}{M} \subseteq \Domain{M}^n$
\item \emph{Interpretasi !!{function}s:} kanggo saben $n$-panggonan
  !!{function}~$f$ ing $\Lang{L}$, fungsi $n$-panggonan
  $\Assign{f}{M} \colon \Domain{M}^n \to \Domain{M}$
\end{enumerate}
\end{defn}

\begin{ex}
!!^a{structure}~$\Struct M$ kanggo basa aritmetika dumadi saka
sawijining himpunan, sawijining elemen ing $\Domain M$,
$\Assign{\Obj 0}{M}$, minangka interpretasi saka !!{constant}~$\Obj 0$,
fungsi siji panggonan $\Assign{\Obj \prime}{M} \colon \Domain{M} \to
\Domain M$, rong fungsi rong panggonan $\Assign{\Obj +}{M}$ lan
$\Assign{\Obj \times}{M}$, kang kalorone $\Domain M^2 \to \Domain M$,
lan relasi rong panggonan $\Assign{\Obj <}{M} \subseteq \Domain{M}^2$.

Tuladha kang cetha saka struktur kaya mangkono yaiku:
\begin{enumerate}
\item $\Domain N = \Nat$
\item $\Assign{\Obj 0}{N} = 0$
\item $\Assign{\Obj \prime}{N}(n) = n + 1$ kanggo saben $n \in \Nat$
\item $\Assign{\Obj +}{N}(n, m) = n + m$ kanggo saben $n, m \in \Nat$
\item $\Assign{\Obj \times}{N}(n, m) = n\cdot m$ kanggo saben $n, m \in \Nat$
\item $\Assign{\Obj <}{N} = \Setabs{\tuple{n, m}}{n \in \Nat, m \in
  \Nat, n < m}$
\end{enumerate}
Struktur~$\Struct N$ kanggo $\Lang L_A$ kang ditetepake mangkono
diarani \emph{model standar aritmetika}, amarga struktur iki menehi
interpretasi marang konstanta nonlogis ing~$\Lang L_A$ persis kaya
kang lumrahe dikarepake.

Nanging, isih ana akeh !!{structure}s liyane kang bisa digawe
kanggo~$\Lang L_A$. Upamane, kita bisa njupuk himpunan~$\Int$
saka wilangan bulat minangka domain, dudu~$\Nat$, banjur nemtokake
interpretasi $\Obj 0$, $\Obj \prime$, $\Obj +$, $\Obj \times$, lan
$\Obj <$ kanthi cocog. Kita uga bisa nemtokake struktur kanggo~$\Lang L_A$
kang babar pisan ora ana gegayutane karo wilangan.
\end{ex}

\begin{ex}
Sawijining struktur~$\Struct M$ kanggo basa~$\Lang L_Z$ saka teori
himpunan mung mbutuhake sawijining himpunan lan siji relasi rong
panggonan. Dadi, sacara teknis, upamane, himpunan wong bebarengan
karo relasi ``$x$ luwih tuwa tinimbang $y$'' bisa digunakake minangka
!!a{structure} kanggo $\Lang L_Z$, semono uga $\Nat$ bebarengan
karo $n \ge m$ kanggo $n, m \in \Nat$.

Salah siji !!{structure} kang mligi narik kawigaten kanggo $\Lang L_Z$,
kanthi !!{element}s ing domaine pancen awujud himpunan lan interpretasi
$\Obj \in$ pancen relasi ``$x$ iku !!a{element} saka~$y$'', yaiku
!!{structure}~$\Struct{{HF}}$ saka \emph{himpunan winates turun-temurun}:
\begin{enumerate}
\item $\Domain{{{HF}}} = \emptyset \cup \Pow{\emptyset} \cup
  \Pow{\Pow{\emptyset}} \cup \Pow{\Pow{\Pow{\emptyset}}} \cup \dots$;
\item $\Assign{\Obj \in}{{{HF}}} = \Setabs{\tuple{x, y}}{x, y \in
  \Domain{{{HF}}}, x \in y}$.
\end{enumerate}
\end{ex}

\begin{digress}
Syarat-syarat kang kita tetepake kanggo nemtokake apa kang dianggep
!!a{structure} nduweni pangaribawa marang logika kita. Upamane,
pilihan supaya domain ora kosong njamin, miturut definisi lumrah
relasi nyukupi (utawa bebener) kanggo ukara mawa kuantor, yen
$\lexists[x][(!A(x) \lor \lnot !A(x))]$ iku valid---yaiku bebener
logis. Lan syarat manawa kabeh !!{constant}s kudu nuduhake obyek
ing domain njamin yen generalisasi eksistensial iku pola inferensi
kang ngemu jaminan semantis: $!A(a)$, mula
$\lexists[x][!A(x)]$. Yen kita ngidini jeneng nuduhake obyek
ing njaba domain, utawa ora nuduhake obyek babar pisan, kita bakal
nyedhaki \emph{logika bebas}, kang mbutuhake premis tambahan kanggo
generalisasi eksistensial: $!A(a)$ lan
$\lexists[x][\eq[x][a]]$, mula $\lexists[x][!A(x)]$.
\end{digress}

\end{document}
